% 01_introduction.tex -- the rehauled paper's section, as bullets to expand into prose (2026-10-15).
% Sources and pointers behind every number: writeup/PAPER_REHAUL_OUTLINE_2026-10-15.md, section 1.
% House rules: every number in the final prose comes from a generated macro (writeup/generated/*.tex), never
% typed; park (comment out with a dated note) rather than delete.
% Macros to use here: \cmNdays, \cmSellPaperMin/Max, \fqsPairs, \fqsQLower(Sig), \fqsShHigher(Sig), \fqsMedDQPct, \fqsMedDSh, \cmLdPerBarBaseQ, \cmLdPerBarBaseDM, \cmMasterRankCorr
% 2026-10-15 (notation pass, writeup/NOTATION_GUIDE_2026-10-15.md): symbols and wording; typed results replaced by
% the macros listed above (values unchanged). Old -> new: RV_t -> \RV_t (closing interval c); B_t -> s^2_{t,c};
% IV^{30}_t -> \mathrm{IV}_t; N -> T_oos; x_t -> \bm x_t; f_{m,d,g} -> f_{m,\mathcal X,\mathcal I} (d clashes with
% the calendar dummies \bm d_t and the loss differential d_t, g with the boosting gradient and the network branch);
% all / HAR+cal -> full predictor set / HAR benchmark; per bar / pooled -> closing-interval / pooled intraday model;
% s_t (signal) -> written out as \hat\sigma^2_t - \mathrm{IV}_t^2 (s is the periodicity factor); Sharpe -> \SR.
% This is the first section to quote numbers, so it loaded their macro files (guarded \input).
% 2026-10-15 (consolidation): calendar dummies \bm d_t -> \bm D_t (d_t is the DM loss differential); the minute
% index k of r_{t,c,k} is a bound summation index only (k is the HAR horizon); upper limit 30 -> N_{t,c}.
% 2026-10-15 (consolidation): parked -- the number files are loaded once in the preamble of paper_rehaul.tex.
% \ifdefined\cmNdays\else% AUTO-GENERATED by writeup/make_table_close_main.py -- do not edit.
% Numbers quoted in the prose of writeup/sections/results_close_option.tex; each is read
% from results/close_master_table/master_table.csv and the sources in the script.
% 2026-10-06: recalibration removed / intervals removed, user decision -- the plain forecast
% f^2 B, point estimates; ...T = HAC t of the daily mid-return difference, ...DM = Diebold-
% Mariano statistic; a ...Above / ...Below count counts point estimates, its ...Sig twin those
% significant at 5 % (DM for QLIKE, HAC t for Sharpe).
\newcommand{\cmNdays}{866}
\newcommand{\cmNrows}{22}
\newcommand{\cmShortMean}{0.014}
\newcommand{\cmShortShMid}{0.20}
\newcommand{\cmShortShX}{\ensuremath{-}0.27}
\newcommand{\cmShortT}{\ensuremath{-}1.71}
\newcommand{\cmPaperShMin}{0.43}
\newcommand{\cmPaperShMax}{1.50}
\newcommand{\cmPaperXMin}{\ensuremath{-}0.05}
\newcommand{\cmPaperXMax}{1.01}
\newcommand{\cmPaperVsShortMin}{+0.22}
\newcommand{\cmPaperVsShortMax}{+1.30}
\newcommand{\cmPaperVsShortAboveN}{8}
\newcommand{\cmPaperVsShortAboveSigN}{three}
\newcommand{\cmPaperVsShortAboveSigWho}{LightGBM, XGBoost, lasso (causally tuned)}
\newcommand{\cmPaperSameMin}{82}
\newcommand{\cmSellPaperMin}{72}
\newcommand{\cmSellPaperMax}{83}
\newcommand{\cmPaperSameMax}{96}
\newcommand{\cmRefQ}{0.1092}
\newcommand{\cmRefShMid}{1.28}
\newcommand{\cmRefShX}{0.80}
\newcommand{\cmRefBuy}{27}
\newcommand{\cmRefVsShort}{+1.08}
\newcommand{\cmRefVsShortT}{+1.71}
\newcommand{\cmFomcSh}{1.36}
\newcommand{\cmFomcD}{+0.08}
\newcommand{\cmFomcDT}{+0.44}
\newcommand{\cmLassoFShMid}{1.06}
\newcommand{\cmLassoFVsRef}{\ensuremath{-}0.22}
\newcommand{\cmLassoFVsRefT}{\ensuremath{-}0.75}
\newcommand{\cmLassoFSame}{96}
\newcommand{\cmPaperFarLabel}{baseline (HAR + calendar OLS)}
\newcommand{\cmPaperFarVsRef}{\ensuremath{-}0.85}
\newcommand{\cmPaperFarVsRefT}{\ensuremath{-}1.49}
\newcommand{\cmPaperBestLabel}{lasso (causally tuned)}
\newcommand{\cmPaperBestSh}{1.50}
\newcommand{\cmNlFamN}{five}
\newcommand{\cmNlMasterN}{90}
\newcommand{\cmPoolShMin}{0.91}
\newcommand{\cmPoolShMax}{1.44}
\newcommand{\cmSubShMin}{1.37}
\newcommand{\cmSubShMax}{1.92}
\newcommand{\cmNlShMin}{1.23}
\newcommand{\cmNlShMax}{1.81}
\newcommand{\cmTabVsRefN}{21}
\newcommand{\cmTabVsRefAbove}{15}
\newcommand{\cmTabVsRefBelow}{6}
\newcommand{\cmTabVsRefAboveSig}{0}
\newcommand{\cmTabVsRefBelowSig}{0}
\newcommand{\cmMasterN}{112}
\newcommand{\cmMasterNPaper}{9}
\newcommand{\cmMasterNBase}{52}
\newcommand{\cmMasterMedQBase}{0.1033}
\newcommand{\cmMasterMedShBase}{0.73}
\newcommand{\cmMasterMedShXBase}{0.25}
\newcommand{\cmMasterNAll}{51}
\newcommand{\cmMasterMedQAll}{0.1014}
\newcommand{\cmMasterMedShAll}{1.27}
\newcommand{\cmMasterMedShXAll}{0.79}
\newcommand{\cmMasterSetsN}{two}
\newcommand{\cmMasterVsRefN}{111}
\newcommand{\cmMasterVsRefAbove}{32}
\newcommand{\cmMasterVsRefBelow}{79}
\newcommand{\cmMasterVsRefAboveSig}{0}
\newcommand{\cmMasterVsRefBelowSig}{0}
\newcommand{\cmMasterVsShortAbove}{112}
\newcommand{\cmMasterVsShortBelow}{0}
\newcommand{\cmMasterVsShortAboveSig}{20}
\newcommand{\cmMasterVsShortBelowSig}{0}
\newcommand{\cmMasterVsShortN}{112}
\newcommand{\cmMasterQVsRefBelow}{99}
\newcommand{\cmMasterQVsRefAbove}{12}
\newcommand{\cmMasterQVsRefBelowSig}{23}
\newcommand{\cmMasterQVsRefAboveSig}{1}
\newcommand{\cmMasterMedQ}{0.1028}
\newcommand{\cmMasterMedSh}{1.08}
\newcommand{\cmMasterMedShX}{0.59}
\newcommand{\cmMasterShMin}{0.29}
\newcommand{\cmMasterShMax}{1.92}
\newcommand{\cmMasterXMin}{\ensuremath{-}0.19}
\newcommand{\cmMasterXMax}{1.43}
\newcommand{\cmMasterVsShortMed}{+0.87}
\newcommand{\cmMasterRankCorr}{\ensuremath{-}0.20}
\newcommand{\cmMasterRankCorrX}{\ensuremath{-}0.20}
\newcommand{\cmSellMasterMin}{66}
\newcommand{\cmSellMasterMax}{83}
\newcommand{\cmSellMasterMed}{70}
\newcommand{\cmTopShLabel}{per-bar elastic net}
\newcommand{\cmTopShInputs}{\texttt{all\_features}}
\newcommand{\cmTopShMid}{1.92}
\newcommand{\cmXLossMin}{0.48}
\newcommand{\cmXLossMax}{0.50}
\newcommand{\cmLdSpec}{+0.16}
\newcommand{\cmLdSpecQ}{\ensuremath{-}4.3}
\newcommand{\cmLdSpecT}{+0.34}
\newcommand{\cmLdSpecDM}{\ensuremath{-}1.73}
\newcommand{\cmLdPerBarBase}{+0.46}
\newcommand{\cmLdPerBarBaseQ}{\ensuremath{-}12.0}
\newcommand{\cmLdPerBarBaseT}{+1.05}
\newcommand{\cmLdPerBarBaseDM}{\ensuremath{-}3.37}
\newcommand{\cmLdPerBarAll}{\ensuremath{-}0.03}
\newcommand{\cmLdPerBarAllQ}{\ensuremath{-}6.7}
\newcommand{\cmLdPerBarAllT}{\ensuremath{-}0.07}
\newcommand{\cmLdPerBarAllDM}{\ensuremath{-}1.61}
\newcommand{\cmLdIn}{+0.05}
\newcommand{\cmLdInQ}{+0.8}
\newcommand{\cmLdInT}{+0.09}
\newcommand{\cmLdInDM}{+0.18}
\newcommand{\cmLdEstMin}{\ensuremath{-}0.15}
\newcommand{\cmLdEstMax}{+0.50}
\newcommand{\cmLdEstMinLabel}{\texttt{subtree\_rf\_all\_features}}
\newcommand{\cmLdEstMaxLabel}{\texttt{subpt\_enet\_all\_features}}
\newcommand{\cmLdEstMinWho}{random forest}
\newcommand{\cmLdEstMaxWho}{elastic net}
\newcommand{\cmLdN}{21}
\newcommand{\cmLdNPath}{12}
\newcommand{\cmLdNCampaign}{9}
\newcommand{\cmLdResolved}{none}
\newcommand{\cmLdQResolved}{three}
\newcommand{\cmLdCDedup}{+0.00}
\newcommand{\cmLdCDedupQ}{+0.0}
\newcommand{\cmLdCMaskTten}{\ensuremath{-}0.21}
\newcommand{\cmLdCMaskTtenQ}{\ensuremath{-}1.0}
\newcommand{\cmLdCMaskTtenT}{\ensuremath{-}0.69}
\newcommand{\cmLdCMaskTtenDM}{\ensuremath{-}1.04}
\newcommand{\cmLdCMaskTone}{\ensuremath{-}0.23}
\newcommand{\cmLdCMaskToneQ}{+1.4}
\newcommand{\cmLdCMaskToneT}{\ensuremath{-}0.55}
\newcommand{\cmLdCMaskToneDM}{+1.36}
\newcommand{\cmLdCMaskRsten}{\ensuremath{-}0.08}
\newcommand{\cmLdCMaskRstenQ}{+2.8}
\newcommand{\cmLdCMaskRstenT}{\ensuremath{-}0.32}
\newcommand{\cmLdCMaskRstenDM}{+1.42}
\newcommand{\cmLdCMaskRsone}{\ensuremath{-}0.05}
\newcommand{\cmLdCMaskRsoneQ}{+0.2}
\newcommand{\cmLdCMaskRsoneT}{\ensuremath{-}0.17}
\newcommand{\cmLdCMaskRsoneDM}{+0.12}
\newcommand{\cmLdCCadLgbm}{\ensuremath{-}0.04}
\newcommand{\cmLdCCadLgbmQ}{\ensuremath{-}2.2}
\newcommand{\cmLdCCadLgbmT}{\ensuremath{-}0.08}
\newcommand{\cmLdCCadLgbmDM}{\ensuremath{-}1.12}
\newcommand{\cmLdCTuneRs}{+0.28}
\newcommand{\cmLdCTuneRsQ}{+3.0}
\newcommand{\cmLdCTuneRsT}{+0.58}
\newcommand{\cmLdCTuneRsDM}{+1.12}
\newcommand{\cmLdCTuneOp}{\ensuremath{-}0.10}
\newcommand{\cmLdCTuneOpQ}{+7.3}
\newcommand{\cmLdCTuneOpT}{\ensuremath{-}0.31}
\newcommand{\cmLdCTuneOpDM}{+2.83}
\newcommand{\cmLdCTuneOpPer}{\ensuremath{-}0.32}
\newcommand{\cmLdCTuneOpPerQ}{+2.5}
\newcommand{\cmLdCTuneOpPerT}{\ensuremath{-}0.65}
\newcommand{\cmLdCTuneOpPerDM}{+0.75}
\newcommand{\cmPNewBase}{22}
\newcommand{\cmPOldBase}{34}
\newcommand{\cmPNewAll}{628}
\newcommand{\cmPOldAll}{640}
\newcommand{\cmNDup}{12}
\newcommand{\cmDedupN}{37}
\newcommand{\cmDedupQMax}{0.05}
\newcommand{\cmDedupPosMax}{one}
\newcommand{\cmKeptBase}{17}
\newcommand{\cmKeptBaseMin}{17}
\newcommand{\cmKeptBaseMax}{17}
\newcommand{\cmKeptAll}{369}
\newcommand{\cmKeptAllMin}{360}
\newcommand{\cmKeptAllMax}{378}
\newcommand{\cmMaskN}{16}
\newcommand{\cmMaskQMed}{0.4}
\newcommand{\cmMaskQAbove}{10}
\newcommand{\cmMaskShBelow}{14}
\newcommand{\cmMaskQAboveSig}{none}
\newcommand{\cmMaskShBelowSig}{one}
\newcommand{\cmTrVsRidgeN}{48}
\newcommand{\cmTrSetsN}{two}
\newcommand{\cmMaskSetsN}{two}
\newcommand{\cmTrVsRidgeQAbove}{46}
\newcommand{\cmTrVsRidgeShBelow}{41}
\newcommand{\cmTrVsRidgeQAboveSig}{21}
\newcommand{\cmTrVsRidgeShBelowSig}{one}
\newcommand{\cmCadN}{six}
\newcommand{\cmCadQBelow}{five}
\newcommand{\cmCadQBelowSig}{none}
\newcommand{\cmCadShUp}{\ensuremath{-}0.43}
\newcommand{\cmCadShUpT}{\ensuremath{-}1.03}
\newcommand{\cmRsQAbove}{6}
\newcommand{\cmRsQAboveSig}{two}
\newcommand{\cmRsN}{6}
\newcommand{\cmOpN}{24}
\newcommand{\cmOpQAbove}{24}
\newcommand{\cmOpQAboveSig}{13}
\newcommand{\cmTpN}{18}
\newcommand{\cmTpQBelow}{12}
\newcommand{\cmTpQBelowSig}{one}
\newcommand{\cmBudN}{24}
\newcommand{\cmBudQAbove}{17}
\newcommand{\cmBudQAboveSig}{two}
\newcommand{\cmBudQBelow}{7}
\newcommand{\cmBudLastMin}{27}
\newcommand{\cmBudLastMax}{37}
\newcommand{\cmBudLastExch}{20}
\newcommand{\cmNlBestSh}{1.81}
\newcommand{\cmNlBestLabel}{tuned per-bar LightGBM}
\newcommand{\cmNlBestInputs}{\texttt{all\_features}}
\newcommand{\cmNlBestVsRef}{+0.53}
\newcommand{\cmNlBestVsRefT}{+1.04}
\newcommand{\cmNlVsRefAbove}{19}
\newcommand{\cmNlVsRefBelow}{71}
\newcommand{\cmNlVsRefAboveSig}{0}
\newcommand{\cmNlVsRefBelowSig}{0}
\newcommand{\cmClassChunks}{505}
\newcommand{\cmClassPts}{20}
\newcommand{\cmClassLgbmPick}{18}
\newcommand{\cmClassMaxRelMin}{4.5}
\newcommand{\cmClassMaxRelMax}{5.5}
\newcommand{\cmSvN}{31}
\newcommand{\cmSvSetN}{112}
\newcommand{\cmSvSetPairsN}{3472}
\newcommand{\cmSvSetAboveMid}{755}
\newcommand{\cmSvSetAboveMidSig}{29}
\newcommand{\cmSvSetStraddleFirstMid}{4}
\newcommand{\cmSvSetWingN}{672}
\newcommand{\cmSvSetWingMidBelow}{668}
\newcommand{\cmSvSetAboveX}{410}
\newcommand{\cmSvSetAboveXSig}{0}
\newcommand{\cmSvSetStraddleFirstX}{3}
\newcommand{\cmSvSetWingXBelow}{672}
\newcommand{\cmSvSetBestMidLabel}{call vertical, w50}
\newcommand{\cmSvSetBestMidMed}{+0.37}
\newcommand{\cmSvSetBestMidAboveN}{94}
\newcommand{\cmSvSetBestMidAboveSig}{0}
\newcommand{\cmSvCostStraddle}{2.9}
\newcommand{\cmSvCostIbfMin}{4.0}
\newcommand{\cmSvCostIbfMax}{6.1}
\newcommand{\cmMasterPosTotalN}{112}
\newcommand{\cmMasterTopTwentyShareMed}{99}
\newcommand{\cmMasterTopTwentyShareMin}{59}
\newcommand{\cmMasterTopTwentyShareMax}{338}
\newcommand{\cmMasterTopTwentyBuysN}{112}
\newcommand{\cmMasterBottomTwentySellsN}{112}
\newcommand{\cmMasterDiffMed}{+53.5}
\newcommand{\cmMasterDiffPosN}{112}
\newcommand{\cmMasterDiffExTenMed}{\ensuremath{-}30.0}
\newcommand{\cmMasterDiffExTenHalfN}{112}
\newcommand{\cmMasterDiffExTenNegN}{102}
\newcommand{\cmBContamMaxDSh}{0.38}
\newcommand{\cmBContamMaxDShX}{0.38}
\newcommand{\cmBContamMedDSh}{0.08}
\newcommand{\cmBContamPaperMaxDSh}{0.21}
\newcommand{\cmBContamPaperMaxDShX}{0.22}
\newcommand{\cmBContamPaperMedDSh}{0.08}
\newcommand{\cmBContamDays}{94}
\newcommand{\cmCensIvMaxDSh}{0.06}
\newcommand{\cmCensIvMaxDShX}{0.08}
\newcommand{\cmCensIvMedDSh}{0.04}
\newcommand{\cmCensIvPaperMaxDSh}{0.05}
\newcommand{\cmCensIvPaperMaxDShX}{0.07}
\newcommand{\cmCensIvPaperMedDSh}{0.04}
\newcommand{\cmCensIvDays}{5}
\newcommand{\cmBContamMaxFlips}{16}
\newcommand{\cmBContamMedFlips}{8}
\newcommand{\cmBContamPaperMaxFlips}{10}
\newcommand{\cmBContamMaxRel}{21}
\newcommand{\cmPtPreSessions}{3947}
\newcommand{\cmPtPreStart}{13~April~2004}
\newcommand{\cmPtValFirst}{6~January~2016}
\newcommand{\cmPtValLast}{2~January~2020}
\newcommand{\cmPtFolds}{four}
\newcommand{\cmPtFoldLen}{250}
\newcommand{\cmPtFitWin}{2{,}000}
\newcommand{\cmPtEmbargo}{25}
\newcommand{\cmPtRatio}{1/100}
\newcommand{\cmPtBlock}{250}
\newcommand{\cmPtRetuneN}{four}
\newcommand{\cmPtRetuneFirst}{23~June~2020}
\newcommand{\cmPtRetuneLast}{16~June~2023}
\newcommand{\cmPtFiredRidgeBase}{one}
\newcommand{\cmPtPenRidgeBase}{\ensuremath{10^{-2}}}
\newcommand{\cmPtFiredRidgeAll}{none}
\newcommand{\cmPtPenRidgeAll}{\ensuremath{10^{3}}}
\newcommand{\cmPtFiredLassoBase}{one}
\newcommand{\cmPtPenLassoBase}{\ensuremath{10^{-6}}}
\newcommand{\cmPtFiredLassoAll}{none}
\newcommand{\cmPtPenLassoAll}{\ensuremath{10^{-2}}}
\newcommand{\cmPtFiredEnetBase}{two}
\newcommand{\cmPtPenEnetBase}{\ensuremath{10^{-6}}}
\newcommand{\cmPtFiredEnetAll}{none}
\newcommand{\cmPtPenEnetAll}{\ensuremath{10^{-2}}}
\newcommand{\cmPtFiredTotal}{4}
\newcommand{\cmPtFiredOf}{24}
\newcommand{\cmPtCausalSwitches}{75}
\newcommand{\cmPtCausalChoices}{165}
\newcommand{\cmPtCausalMarginPass}{3}
\newcommand{\cmPtCausalEdgePicks}{132}
\newcommand{\cmPtCausalPicks}{198}
\newcommand{\cmPtSwitchFlipPct}{8}
\newcommand{\cmPtCausalSwitchesPaper}{24}
\newcommand{\cmPtCausalChoicesPaper}{45}
\newcommand{\cmPtCausalMarginPassPaper}{3}
\newcommand{\cmPtCausalEdgePicksPaper}{34}
\newcommand{\cmPtCausalPicksPaper}{54}
\newcommand{\cmPtSwitchFlipPctPaper}{9}
\newcommand{\cmPoolPtN}{6}
\newcommand{\cmPoolPtQPctMean}{+0.9}
\newcommand{\cmPoolPtQWorse}{5}
\newcommand{\cmPoolPtDMMax}{1.22}
\newcommand{\cmPoolPtTMax}{1.29}
\newcommand{\cmPoolPtSigN}{0}
\newcommand{\cmPoolPtShAllMin}{+0.17}
\newcommand{\cmPoolPtShAllMax}{+0.35}
\newcommand{\cmPoolPtShOtherMin}{\ensuremath{-}0.06}
\newcommand{\cmPoolPtShOtherMax}{\ensuremath{-}0.19}
\newcommand{\cmPoolPtBlocks}{218}
\newcommand{\cmPoolPtCausalSwMin}{137}
\newcommand{\cmPoolPtCausalSwMax}{170}
\newcommand{\cmPoolPtCausalMarginMin}{3}
\newcommand{\cmPoolPtCausalMarginMax}{22}
\newcommand{\cmPoolPtRetunes}{217}
\newcommand{\cmPoolPtRetuneSwMin}{11}
\newcommand{\cmPoolPtRetuneSwMax}{40}
\newcommand{\cmTreeProvRfName}{\texttt{rf\_opt\_00}}
\newcommand{\cmTreeProvRfTrials}{400}
\newcommand{\cmTreeProvRfBudget}{400}
\newcommand{\cmTreeProvRfDevMse}{0.108456}
\newcommand{\cmTreeProvLgbmDevMse}{0.107449}
\newcommand{\cmTreeProvXgbDevMse}{0.108076}
\newcommand{\cmTreeProvXgbTopDevMse}{0.107894}
\newcommand{\cmTreeProvRfTrees}{485}
\newcommand{\cmTreeProvRfDepth}{20}
\newcommand{\cmTreeProvRfFeatures}{0.90}
\newcommand{\cmTreeProvRfSamples}{0.41}
\newcommand{\cmTreeProvRfLeafDev}{12}
\newcommand{\cmTreeProvRfLeaf}{1}
\newcommand{\cmTreeProvLgbmLeafDev}{98}
\newcommand{\cmTreeProvLgbmLeaf}{8}
\newcommand{\cmTreeProvXgbTopName}{\texttt{xgb\_opt\_00}}
\newcommand{\cmTreeProvXgbName}{\texttt{xgb\_opt\_06}}
\newcommand{\cmTreeProvXgbDevRank}{7}
\newcommand{\cmTreeProvXgbArms}{8}
\newcommand{\cmTreeProvXgbPostRank}{8}
\newcommand{\cmTreeProvXgbRecentRank}{8}
\newcommand{\cmTreeProvXgbQRank}{8}
\newcommand{\cmTreeProvXgbPickDate}{14~August~2026}
\newcommand{\cmMasterNLean}{51}
\newcommand{\cmMasterMedQLean}{0.0991}
\newcommand{\cmMasterMedShLean}{1.37}
\newcommand{\cmMasterMedShXLean}{0.88}
\newcommand{\cmLdPerBarLean}{+0.65}
\newcommand{\cmLdPerBarLeanQ}{\ensuremath{-}12.3}
\newcommand{\cmLdPerBarLeanT}{+1.49}
\newcommand{\cmLdPerBarLeanDM}{\ensuremath{-}3.39}
\newcommand{\cmLdInBaseLean}{+0.51}
\newcommand{\cmLdInBaseLeanQ}{\ensuremath{-}3.3}
\newcommand{\cmLdInBaseLeanT}{+1.66}
\newcommand{\cmLdInBaseLeanDM}{\ensuremath{-}1.17}
\newcommand{\cmLdInLeanAll}{\ensuremath{-}0.46}
\newcommand{\cmLdInLeanAllQ}{+4.3}
\newcommand{\cmLdInLeanAllT}{\ensuremath{-}1.15}
\newcommand{\cmLdInLeanAllDM}{+1.50}
%
% 2026-10-06 parked (recalibration removed / intervals removed, user decision):
%   \cmShortShCiMid: interval of always short's Sharpe
%   \cmShortShCiX: interval of always short's Sharpe
%   \cmPaperVsShortAboveWho: the paper's forecasts whose interval vs always short lay above zero
%   \cmPaperVsShortAboveCi: that forecast's interval
%   \cmSellHead: the former headline's sell share (no headline model)
%   \cmHQ: the former headline (no headline model)
%   \cmHQPct: the former headline (no headline model)
%   \cmHQCi: the former headline (no headline model)
%   \cmHDM: the former headline (no headline model)
%   \cmHDMp: the former headline (no headline model)
%   \cmHShMid: the former headline (no headline model)
%   \cmHShX: the former headline (no headline model)
%   \cmHMean: the former headline (no headline model)
%   \cmHHit: the former headline (no headline model)
%   \cmHBuy: the former headline (no headline model)
%   \cmHVsRefMid: the former headline (no headline model)
%   \cmHVsRefX: the former headline (no headline model)
%   \cmHVsShortMid: the former headline (no headline model)
%   \cmHVsShortX: the former headline (no headline model)
%   \cmHVsRefT: the former headline (no headline model)
%   \cmHVsRefP: the former headline (no headline model)
%   \cmHVsShortT: the former headline (no headline model)
%   \cmHVsShortP: the former headline (no headline model)
%   \cmHRankSh: the former headline (no headline model)
%   \cmHRankQ: the former headline (no headline model)
%   \cmVsHN: the former headline (no headline model)
%   \cmVsHAbove: the former headline (no headline model)
%   \cmVsHBelow: the former headline (no headline model)
%   \cmVsHAboveSig: the former headline (no headline model)
%   \cmVsHBelowSig: the former headline (no headline model)
%   \cmVsHLassoSh: the former headline (no headline model)
%   \cmVsHLasso: paired against the former headline (no headline model)
%   \cmVsHLassoT: paired against the former headline (no headline model)
%   \cmVsHLassoSame: paired against the former headline (no headline model)
%   \cmVsHLassoQPct: paired against the former headline (no headline model)
%   \cmVsHLassoDM: paired against the former headline (no headline model)
%   \cmVsHEnet: paired against the former headline (no headline model)
%   \cmVsHEnetT: paired against the former headline (no headline model)
%   \cmVsHEnetSame: paired against the former headline (no headline model)
%   \cmVsHEnetQPct: paired against the former headline (no headline model)
%   \cmVsHEnetDM: paired against the former headline (no headline model)
%   \cmVsHRidgeAll: paired against the former headline (no headline model)
%   \cmVsHRidgeAllT: paired against the former headline (no headline model)
%   \cmVsHRidgeAllSame: paired against the former headline (no headline model)
%   \cmVsHRidgeAllQPct: paired against the former headline (no headline model)
%   \cmVsHRidgeAllDM: paired against the former headline (no headline model)
%   \cmVsHTop: paired against the former headline (no headline model)
%   \cmVsHTopT: paired against the former headline (no headline model)
%   \cmVsHTopSame: paired against the former headline (no headline model)
%   \cmVsHTopQPct: paired against the former headline (no headline model)
%   \cmVsHTopDM: paired against the former headline (no headline model)
%   \cmLdPerBarLive: a step to the former headline (no headline, no live-feasible)
%   \cmLdPerBarLiveQ: a step to the former headline (no headline, no live-feasible)
%   \cmLdPerBarLiveT: a step to the former headline (no headline, no live-feasible)
%   \cmLdPerBarLiveDM: a step to the former headline (no headline, no live-feasible)
%   \cmLdInAll: a step to the former headline (no headline, no live-feasible)
%   \cmLdInAllQ: a step to the former headline (no headline, no live-feasible)
%   \cmLdInAllT: a step to the former headline (no headline, no live-feasible)
%   \cmLdInAllDM: a step to the former headline (no headline, no live-feasible)
%   \cmLdInBase: a step to the former headline (no headline, no live-feasible)
%   \cmLdInBaseQ: a step to the former headline (no headline, no live-feasible)
%   \cmLdInBaseT: a step to the former headline (no headline, no live-feasible)
%   \cmLdInBaseDM: a step to the former headline (no headline, no live-feasible)
%   \cmLdCMaskLstm: an LSTM step of the campaign (LSTMs parked)
%   \cmLdCMaskLstmQ: an LSTM step of the campaign (LSTMs parked)
%   \cmLdCMaskLstmT: an LSTM step of the campaign (LSTMs parked)
%   \cmLdCMaskLstmDM: an LSTM step of the campaign (LSTMs parked)
%   \cmLdCCadLstm: an LSTM step of the campaign (LSTMs parked)
%   \cmLdCCadLstmQ: an LSTM step of the campaign (LSTMs parked)
%   \cmLdCCadLstmT: an LSTM step of the campaign (LSTMs parked)
%   \cmLdCCadLstmDM: an LSTM step of the campaign (LSTMs parked)
%   \cmLdCSeq: an LSTM step of the campaign (LSTMs parked)
%   \cmLdCSeqQ: an LSTM step of the campaign (LSTMs parked)
%   \cmLdCSeqT: an LSTM step of the campaign (LSTMs parked)
%   \cmLdCSeqDM: an LSTM step of the campaign (LSTMs parked)
%   \cmPNewLive: live_feasible input set (live-feasible removed from the paper)
%   \cmPOldLive: live_feasible input set (live-feasible removed from the paper)
%   \cmKeptLive: live_feasible input set (live-feasible removed from the paper)
%   \cmKeptLiveMin: live_feasible input set (live-feasible removed from the paper)
%   \cmKeptLiveMax: live_feasible input set (live-feasible removed from the paper)
%   \cmLstmCadLive: live_feasible input set (live-feasible removed from the paper)
%   \cmLstmCadBase: LSTM (LSTMs parked)
%   \cmLstmCadAll: LSTM (LSTMs parked)
%   \cmLstmCadBelow: LSTM (LSTMs parked)
%   \cmLstmCadBelowSig: LSTM (LSTMs parked)
%   \cmLstmVsRidgeMin: LSTM (LSTMs parked)
%   \cmLstmVsRidgeMax: LSTM (LSTMs parked)
%   \cmIntraVsRidgeMin: LSTM (LSTMs parked)
%   \cmIntraVsRidgeMax: LSTM (LSTMs parked)
%   \cmIntraAllQ: LSTM (LSTMs parked)
%   \cmIntraAllDM: LSTM (LSTMs parked)
%   \cmNlBestVsH: against the former headline (no headline model)
%   \cmNlBestVsHT: against the former headline (no headline model)
%   \cmNlVsHAbove: against the former headline (no headline model)
%   \cmNlVsHAboveSig: against the former headline (no headline model)
%   \cmPnlTotal: D.14, the former headline's P&L decomposition (parked, user decision)
%   \cmPnlTotalX: D.14, the former headline's P&L decomposition (parked, user decision)
%   \cmPnlTopTenShare: D.14, the former headline's P&L decomposition (parked, user decision)
%   \cmPnlTopTwentyShare: D.14, the former headline's P&L decomposition (parked, user decision)
%   \cmPnlTopTwentyShareX: D.14, the former headline's P&L decomposition (parked, user decision)
%   \cmPnlExTwenty: D.14, the former headline's P&L decomposition (parked, user decision)
%   \cmPnlExTwentySh: D.14, the former headline's P&L decomposition (parked, user decision)
%   \cmPnlMktTwenty: D.14, the former headline's P&L decomposition (parked, user decision)
%   \cmPnlMktTwentyExp: D.14, the former headline's P&L decomposition (parked, user decision)
%   \cmPnlMktTwentyP: D.14, the former headline's P&L decomposition (parked, user decision)
%   \cmPnlDaysAll: D.14, the former headline's P&L decomposition (parked, user decision)
%   \cmPnlDaysAllPct: D.14, the former headline's P&L decomposition (parked, user decision)
%   \cmPnlMEn: D.14, the former headline's P&L decomposition (parked, user decision)
%   \cmPnlMEbuys: D.14, the former headline's P&L decomposition (parked, user decision)
%   \cmPnlME: D.14, the former headline's P&L decomposition (parked, user decision)
%   \cmPnlMEShort: D.14, the former headline's P&L decomposition (parked, user decision)
%   \cmPnlNotMEn: D.14, the former headline's P&L decomposition (parked, user decision)
%   \cmPnlNotME: D.14, the former headline's P&L decomposition (parked, user decision)
%   \cmPnlDiffBuyShare: D.14, the former headline's P&L decomposition (parked, user decision)
%   \cmPnlDiff: D.14, the former headline's P&L decomposition (parked, user decision)
%   \cmPnlDiffX: D.14, the former headline's P&L decomposition (parked, user decision)
%   \cmPnlDiffExTen: D.14, the former headline's P&L decomposition (parked, user decision)
%   \cmPnlDiffHighVix: D.14, the former headline's P&L decomposition (parked, user decision)
%   \cmPnlHighVixN: D.14, the former headline's P&L decomposition (parked, user decision)
%   \cmSvAbove: strategy variations under the former headline (no headline model)
%   \cmSvAboveSig: strategy variations under the former headline (no headline model)
%   \cmSvBestMid: strategy variations under the former headline (no headline model)
%   \cmSvBestMidSh: strategy variations under the former headline (no headline model)
%   \cmSvBestMidT: strategy variations under the former headline (no headline model)
%   \cmSvWingN: strategy variations under the former headline (no headline model)
%   \cmSvWingMid: strategy variations under the former headline (no headline model)
%   \cmSvWingX: strategy variations under the former headline (no headline model)
%   \cmSvWingMidSig: strategy variations under the former headline (no headline model)
%   \cmSvWingXSig: strategy variations under the former headline (no headline model)
%   \cmPtFiredRidgeLean: 2026-10-08 parked (lean set; replaced by an MCS)
%   \cmPtPenRidgeLean: 2026-10-08 parked (lean set; replaced by an MCS)
%   \cmPtFiredLassoLean: 2026-10-08 parked (lean set; replaced by an MCS)
%   \cmPtPenLassoLean: 2026-10-08 parked (lean set; replaced by an MCS)
%   \cmPtFiredEnetLean: 2026-10-08 parked (lean set; replaced by an MCS)
%   \cmPtPenEnetLean: 2026-10-08 parked (lean set; replaced by an MCS)
%
% statements of the text checked on these numbers (2026-10-06; not asserted):
%   holds   5.4: sign(s) lifts every one of the paper's eight forecasts above always short (mid)
%   holds   5.4: R trades above always short (cmRefVsShort > 0)
%   holds   5.4: the FOMC columns' Sharpe gap is not significant (HAC t)
%   holds   4.5: every forecast of the model set leans short (sells on > half the days)
%   holds   5.4: the model set's top Sharpe (subpt_enet_all_features) is not significantly above R (HAC t)
%   holds   5.4: no step of Table 8 moves the Sharpe ratio significantly (HAC t; 0 do)
%   holds   5.4: daily refits make LightGBM more accurate (cmLdCCadLgbmQ < 0)
%   holds   5.4: ... without moving its trade significantly (HAC t)
%   holds   5.4: tuning does not help (TuneRs: QLIKE % > 0)
%   holds   5.4: tuning does not help (TuneOp: QLIKE % > 0)
%   holds   5.4: the de-dup moves the per-bar ridge's QLIKE only at rounding
%   holds   5.4: the de-dup leaves the trade
%   holds   5.4: no position of the per-bar ridge changed
%   FLIPPED 5.4: the specification step moves little on the trade (cmLdSpec)
%   holds   5.4: per-bar beats pooled on QLIKE on HAR + calendar (cmLdPerBarBaseQ < 0)
%   holds   5.4: ... significantly (DM)
%   holds   5.4: ... but not on the trade (HAC t not significant)
%   holds   5.4: no linear forecast's trade moved significantly with the de-dup (HAC t)
%   holds   5.4: the mask made no tree significantly better (DM / HAC t)
%   FLIPPED 5.4: no tree more accurate than the ridge (2 QLIKE differences below 0)
%   FLIPPED 5.4: no tree trades above the ridge (7 Sharpe differences above 0)
%   FLIPPED 5.4: daily refits lower every tree's QLIKE
%   holds   5.4: daily refits move XGBoost on all_features' trade the most
%   FLIPPED 5.4 (was): that move is resolved (HAC t significant)
%   holds   5.4: random search is never more accurate than the shipped settings (0 below)
%   holds   5.4: Optuna is never more accurate than the shipped settings (0 below)
%   holds   5.4: no tuned tree's trade differs significantly from the shipped one (0)
%   FLIPPED 5.4: tuning more often is never less accurate (6 above)
%   holds   5.4: ... and never significantly different on the trade (HAC t)
%   FLIPPED 5.4: the best of 50 trials is never more accurate than of 10 / 25 (7 below)
%   FLIPPED 5.4: at crossed fills the straddle ranks first for every forecast of the model set
%   holds   5.4: wings cost at every width at crossed fills, for every forecast
%   holds   5.4: for every forecast the 20 best days are buys (the tail-trade reading)
%   FLIPPED 4.2: the B contamination moves no Sharpe ratio by 0.1 (model set)
%   FLIPPED 4.5: dropping the censored days moves no Sharpe ratio by more than 0.06
%   holds   4.x: the margin rule lets few retunes fire
%   FLIPPED methods: pre-tuning improves the pooled models' QLIKE (mean change < 0)
\fi
% \ifdefined\fqsPairs\else% AUTO-GENERATED by writeup/make_fig_qlike_sharpe.py -- do not edit.
% Figure writeup/figures/fig_qlike_sharpe.pdf: QLIKE against Sharpe (mid), HAR + calendar vs all
% features, every model of results/close_master_table/paper_model_set.csv on both input sets
% (per-bar linear = the pre-tuned subpt_* rows; LSTMs excluded); numbers from master_table.csv
% and master_table_daily.parquet.  Changes are all features minus HAR + calendar, on the same days.
% DM: dm_test on the daily QLIKE difference; t: hac_inference.llsw_test on the daily mid-return
% difference (614 post-warmup days, bandwidth 33, fixed-b p); Sharpe changes: Ledoit-Wolf p;
% ES(c = 1): close_murphy_splits_headline.elementary on f / iv, y / iv.  Point estimates.
% pooled pairs: family pooled twin
\newcommand{\fqsPairs}{42}  % models on both input sets (pairs drawn)
\newcommand{\fqsPairsPerBarLinear}{3}  % of which per-bar linear (pre-tuned)
\newcommand{\fqsPairsPooled}{3}  % of which pooled 48-bar linear
\newcommand{\fqsPairsTree}{36}  % of which per-bar trees
\newcommand{\fqsPoints}{85}  % points drawn: both ends of every pair + the HAR + cal. OLS
\newcommand{\fqsDays}{866}  % trade days
\newcommand{\fqsQLower}{31}  % QLIKE lower with all features
\newcommand{\fqsQLowerSig}{3}  % ... and DM significant at 5 %
\newcommand{\fqsQHigher}{11}  % QLIKE higher with all features
\newcommand{\fqsQHigherSig}{0}  % ... and DM significant at 5 %
\newcommand{\fqsQLowerSigLagFourteen}{2}  % sensitivity: QLIKE lower, DM at lag 14 significant
\newcommand{\fqsQHigherSigLagFourteen}{0}  % sensitivity: QLIKE higher, DM at lag 14 significant
\newcommand{\fqsShHigher}{41}  % Sharpe (mid) higher with all features
\newcommand{\fqsShHigherSig}{6}  % ... and Ledoit-Wolf p < 0.05 (614 post-warmup days)
\newcommand{\fqsShLower}{1}  % Sharpe (mid) lower with all features
\newcommand{\fqsShLowerSig}{0}  % ... and Ledoit-Wolf p < 0.05 (614 post-warmup days)
\newcommand{\fqsQLowerShHigher}{30}  % both: QLIKE lower and Sharpe higher
\newcommand{\fqsMedDQ}{\ensuremath{-}0.0016}  % median QLIKE change (all features minus HAR + cal.)
\newcommand{\fqsMedDQPct}{\ensuremath{-}1.6}  % median QLIKE change, % of the HAR + cal. QLIKE
\newcommand{\fqsMedDSh}{\ensuremath{+}0.53}  % median Sharpe (mid) change
\newcommand{\fqsESLower}{29}  % elementary score at the threshold (c = 1) lower with all features
\newcommand{\fqsCorrES}{0.07}  % Pearson corr. across pairs: ES(c = 1) change vs HAC t of the P&L change
\newcommand{\fqsCorrQ}{0.12}  % Pearson corr. across pairs: QLIKE change vs HAC t of the P&L change
\newcommand{\fqsInShade}{4}  % all-features points in the shaded quadrant
\newcommand{\fqsInShadeBase}{3}  % HAR + cal. points in the shaded quadrant
\newcommand{\fqsOlsQ}{0.0967}  % QLIKE of the HAR + cal. OLS
\newcommand{\fqsOlsSh}{1.06}  % Sharpe (mid) of the HAR + cal. OLS
\newcommand{\fqsDMLag}{6}  % Newey-West lag of the DM statistic (the master table's)
\newcommand{\fqsHACS}{33}  % bandwidth S of the t of the mean P&L change (LLSW 2018, 614 post-warmup days)
\newcommand{\fqsSigPct}{5}  % two-sided level of the counts, %
\newcommand{\fqsQuadBetterBoth}{30}  % change figure: lower QLIKE, higher Sharpe
\newcommand{\fqsQuadShOnly}{11}  % change figure: higher QLIKE, higher Sharpe
\newcommand{\fqsQuadQOnly}{1}  % change figure: lower QLIKE, lower Sharpe
\newcommand{\fqsQuadWorseBoth}{0}  % change figure: higher QLIKE, lower Sharpe
\fi
% \ifdefined\rhfiNExog\else% AUTO-GENERATED by writeup/make_rehaul_fi.py -- do not edit; re-run the script.
% \rhfi... macros for writeup/rehaul/08_feature_importance.tex (results/rehaul_fi/).
\providecommand{\rhfiNRefits}{}\renewcommand{\rhfiNRefits}{147}  % refits, one every 10 sessions
\providecommand{\rhfiNOos}{}\renewcommand{\rhfiNOos}{1{,}469}  % scored held-out sessions
\providecommand{\rhfiP}{}\renewcommand{\rhfiP}{628}  % columns, all features
\providecommand{\rhfiNUnits}{}\renewcommand{\rhfiNUnits}{49}  % units: 47 exogenous series + HAR ladder + calendar
\providecommand{\rhfiNExog}{}\renewcommand{\rhfiNExog}{47}  % exogenous series
\providecommand{\rhfiFirst}{}\renewcommand{\rhfiFirst}{25~June~2018}  % first scored session
\providecommand{\rhfiLast}{}\renewcommand{\rhfiLast}{30~April~2024}  % last scored session
\providecommand{\rhfiWindow}{}\renewcommand{\rhfiWindow}{2{,}000}  % sessions in a training window
\providecommand{\rhfiRefitEvery}{}\renewcommand{\rhfiRefitEvery}{10}  % sessions between refits
\providecommand{\rhfiBackbonePf}{}\renewcommand{\rhfiBackbonePf}{0.01}  % backbone penalty factor
\providecommand{\rhfiNBuckets}{}\renewcommand{\rhfiNBuckets}{8}  % exogenous buckets
\providecommand{\rhfiOldBlock}{}\renewcommand{\rhfiOldBlock}{250}  % sessions between the old arm's penalty re-choices
\providecommand{\rhfiNPairs}{}\renewcommand{\rhfiNPairs}{1{,}176}  % series pairs
\providecommand{\rhfiNProbeCols}{}\renewcommand{\rhfiNProbeCols}{631}  % columns of a probe fit
\providecommand{\rhfiTopWPairs}{}\renewcommand{\rhfiTopWPairs}{50}  % pairs in the interaction W
\providecommand{\rhfiProbeLevels}{}\renewcommand{\rhfiProbeLevels}{20}  % values of the few-valued noise column
\providecommand{\rhfiRidgeAlpha}{}\renewcommand{\rhfiRidgeAlpha}{\ensuremath{10^{3}}}
\providecommand{\rhfiLassoAlpha}{}\renewcommand{\rhfiLassoAlpha}{\ensuremath{10^{-2}}}
\providecommand{\rhfiLassoNzMed}{}\renewcommand{\rhfiLassoNzMed}{30}  % median nonzero coefficients per refit
\providecommand{\rhfiLassoNzMin}{}\renewcommand{\rhfiLassoNzMin}{27}
\providecommand{\rhfiLassoNzMax}{}\renewcommand{\rhfiLassoNzMax}{35}
\providecommand{\rhfiLassoExogNzMed}{}\renewcommand{\rhfiLassoExogNzMed}{14}  % ... exogenous columns
\providecommand{\rhfiLassoColsEver}{}\renewcommand{\rhfiLassoColsEver}{60}  % columns nonzero in >= 1 refit
\providecommand{\rhfiLassoColsMajority}{}\renewcommand{\rhfiLassoColsMajority}{29}  % columns nonzero in >= half the refits
\providecommand{\rhfiLassoExogEver}{}\renewcommand{\rhfiLassoExogEver}{27}  % exogenous series with a nonzero column in >= 1 refit
\providecommand{\rhfiLassoExogMajority}{}\renewcommand{\rhfiLassoExogMajority}{11}  % ... in >= half the refits
\providecommand{\rhfiLassoExogAlways}{}\renewcommand{\rhfiLassoExogAlways}{3}  % ... in every refit
\providecommand{\rhfiLassoAlwaysList}{}\renewcommand{\rhfiLassoAlwaysList}{\texttt{sumret}, \texttt{sumabsret\_vwstock}, \texttt{sumautocov}}  % exogenous series selected in every refit
\providecommand{\rhfiLassoProbeSel}{}\renewcommand{\rhfiLassoProbeSel}{1}  % continuous probe selected, of the probe fits
\providecommand{\rhfiLassoProbeSelLevels}{}\renewcommand{\rhfiLassoProbeSelLevels}{0}
\providecommand{\rhfiLassoProbeSelBinary}{}\renewcommand{\rhfiLassoProbeSelBinary}{0}
\providecommand{\rhfiNProbe}{}\renewcommand{\rhfiNProbe}{49}  % probe refits (every third)
\providecommand{\rhfiLassoW}{}\renewcommand{\rhfiLassoW}{0.74}  % Kendall W of the per-refit series ranks by sum |beta|
\providecommand{\rhfiLassoBackboneShare}{}\renewcommand{\rhfiLassoBackboneShare}{88}  % median share of sum |beta| on the backbone, %
\providecommand{\rhfiOldLassoSpearman}{}\renewcommand{\rhfiOldLassoSpearman}{0.70}  % Spearman, selection share: pre-tuned vs re-tuned one-penalty lasso
\providecommand{\rhfiOldLassoNzMed}{}\renewcommand{\rhfiOldLassoNzMed}{23}  % re-tuned lasso: median nonzero per refit
\providecommand{\rhfiOldLassoExogMajority}{}\renewcommand{\rhfiOldLassoExogMajority}{18}
\providecommand{\rhfiRidgeTopOne}{}\renewcommand{\rhfiRidgeTopOne}{Calendar}
\providecommand{\rhfiRidgeTopOneShare}{}\renewcommand{\rhfiRidgeTopOneShare}{31.0}  % % of sum |beta|
\providecommand{\rhfiRidgeTopTwo}{}\renewcommand{\rhfiRidgeTopTwo}{HAR ladder}
\providecommand{\rhfiRidgeTopTwoShare}{}\renewcommand{\rhfiRidgeTopTwoShare}{22.8}  % % of sum |beta|
\providecommand{\rhfiRidgeTopThree}{}\renewcommand{\rhfiRidgeTopThree}{\texttt{sumret}}
\providecommand{\rhfiRidgeTopThreeShare}{}\renewcommand{\rhfiRidgeTopThreeShare}{3.6}  % % of sum |beta|
\providecommand{\rhfiRidgeTopFour}{}\renewcommand{\rhfiRidgeTopFour}{\texttt{sumabsret\_vwstock}}
\providecommand{\rhfiRidgeTopFourShare}{}\renewcommand{\rhfiRidgeTopFourShare}{2.2}  % % of sum |beta|
\providecommand{\rhfiRidgeSdTopOne}{}\renewcommand{\rhfiRidgeSdTopOne}{HAR ladder}
\providecommand{\rhfiRidgeSdTopOneShare}{}\renewcommand{\rhfiRidgeSdTopOneShare}{30.6}  % % of sum |beta| sd
\providecommand{\rhfiRidgeSdTopTwo}{}\renewcommand{\rhfiRidgeSdTopTwo}{Calendar}
\providecommand{\rhfiRidgeSdTopTwoShare}{}\renewcommand{\rhfiRidgeSdTopTwoShare}{9.3}  % % of sum |beta| sd
\providecommand{\rhfiRidgeSdTopThree}{}\renewcommand{\rhfiRidgeSdTopThree}{\texttt{sumret}}
\providecommand{\rhfiRidgeSdTopThreeShare}{}\renewcommand{\rhfiRidgeSdTopThreeShare}{4.5}  % % of sum |beta| sd
\providecommand{\rhfiRidgeTopExog}{}\renewcommand{\rhfiRidgeTopExog}{\texttt{sumret}}
\providecommand{\rhfiRidgeTopExogShare}{}\renewcommand{\rhfiRidgeTopExogShare}{3.6}
\providecommand{\rhfiRidgeBackboneShare}{}\renewcommand{\rhfiRidgeBackboneShare}{54}  % median share of sum |beta| on the backbone, %
\providecommand{\rhfiRidgeW}{}\renewcommand{\rhfiRidgeW}{0.85}
\providecommand{\rhfiRidgeAllSelected}{}\renewcommand{\rhfiRidgeAllSelected}{49}  % units with N_A = K under the ridge
\providecommand{\rhfiLassoProbeRate}{}\renewcommand{\rhfiLassoProbeRate}{2.0}  % continuous probe's selection rate, %
\providecommand{\rhfiLassoColsAboveProbe}{}\renewcommand{\rhfiLassoColsAboveProbe}{54}  % columns selected more often than the continuous probe
\providecommand{\rhfiLassoSeriesAboveProbe}{}\renewcommand{\rhfiLassoSeriesAboveProbe}{23}  % exogenous series with such a column
\providecommand{\rhfiLassoTopCol}{}\renewcommand{\rhfiLassoTopCol}{\texttt{is\_month\_end}}
\providecommand{\rhfiRidgeSpearmanSd}{}\renewcommand{\rhfiRidgeSpearmanSd}{0.85}  % Spearman, series sum|beta| vs sum|beta| sd
\providecommand{\rhfiRidgeProbeBeat}{}\renewcommand{\rhfiRidgeProbeBeat}{107}  % columns with |beta| above the continuous probe's in most probe fits
\providecommand{\rhfiRidgeProbeBeatSeries}{}\renewcommand{\rhfiRidgeProbeBeatSeries}{38}
\providecommand{\rhfiScaleValueN}{}\renewcommand{\rhfiScaleValueN}{282}
\providecommand{\rhfiScaleValueSd}{}\renewcommand{\rhfiScaleValueSd}{0.84}  % median window sd, value ladders
\providecommand{\rhfiScaleValueSdLo}{}\renewcommand{\rhfiScaleValueSdLo}{0.03}
\providecommand{\rhfiScaleValueSdHi}{}\renewcommand{\rhfiScaleValueSdHi}{1.76}
\providecommand{\rhfiScaleValueIqr}{}\renewcommand{\rhfiScaleValueIqr}{0.94}  % median window IQR, value ladders
\providecommand{\rhfiScaleFlagN}{}\renewcommand{\rhfiScaleFlagN}{324}
\providecommand{\rhfiScaleFlagSd}{}\renewcommand{\rhfiScaleFlagSd}{0.022}  % median window sd, availability / activity ladders
\providecommand{\rhfiScaleFlagSdHi}{}\renewcommand{\rhfiScaleFlagSdHi}{0.17}
\providecommand{\rhfiScaleCalSd}{}\renewcommand{\rhfiScaleCalSd}{0.39}  % median window sd, calendar
\providecommand{\rhfiScaleHarSd}{}\renewcommand{\rhfiScaleHarSd}{1.07}  % median window sd, HAR ladder
\providecommand{\rhfiLgbmW}{}\renewcommand{\rhfiLgbmW}{0.57}  % lgbm: Kendall W, per-refit series ranks
\providecommand{\rhfiLgbmWCols}{}\renewcommand{\rhfiLgbmWCols}{0.91}  % lgbm: Kendall W, column ranks
\providecommand{\rhfiLgbmTopOne}{}\renewcommand{\rhfiLgbmTopOne}{HAR ladder}
\providecommand{\rhfiLgbmTopOneRank}{}\renewcommand{\rhfiLgbmTopOneRank}{1.0}  % mean rank
\providecommand{\rhfiLgbmTopTwo}{}\renewcommand{\rhfiLgbmTopTwo}{\texttt{sumret}}
\providecommand{\rhfiLgbmTopTwoRank}{}\renewcommand{\rhfiLgbmTopTwoRank}{2.9}  % mean rank
\providecommand{\rhfiLgbmTopThree}{}\renewcommand{\rhfiLgbmTopThree}{\texttt{sumabsret}}
\providecommand{\rhfiLgbmTopThreeRank}{}\renewcommand{\rhfiLgbmTopThreeRank}{4.8}  % mean rank
\providecommand{\rhfiLgbmSpBorda}{}\renewcommand{\rhfiLgbmSpBorda}{0.90}  % Spearman, mean rank vs top-10 Borda
\providecommand{\rhfiLgbmSpGlobal}{}\renewcommand{\rhfiLgbmSpGlobal}{0.96}  % Spearman, mean rank vs global mean |phi|
\providecommand{\rhfiLgbmSpLocal}{}\renewcommand{\rhfiLgbmSpLocal}{1.000}  % Spearman, mean ranks local refit vs stored
\providecommand{\rhfiLgbmProbeRank}{}\renewcommand{\rhfiLgbmProbeRank}{83}  % median rank of the continuous probe, 631 columns
\providecommand{\rhfiLgbmProbeCols}{}\renewcommand{\rhfiLgbmProbeCols}{60}  % columns above the continuous probe in most probe fits
\providecommand{\rhfiLgbmProbeSeries}{}\renewcommand{\rhfiLgbmProbeSeries}{27}  % exogenous series with such a column
\providecommand{\rhfiLgbmTopCol}{}\renewcommand{\rhfiLgbmTopCol}{\texttt{har\_ma\_1}}  % column with the lowest mean rank
\providecommand{\rhfiLgbmTopColRank}{}\renewcommand{\rhfiLgbmTopColRank}{1.3}
\providecommand{\rhfiLgbmTopColTwo}{}\renewcommand{\rhfiLgbmTopColTwo}{\texttt{har\_ma\_5}}
\providecommand{\rhfiLgbmTopColTwoRank}{}\renewcommand{\rhfiLgbmTopColTwoRank}{1.8}
\providecommand{\rhfiLgbmPredGap}{}\renewcommand{\rhfiLgbmPredGap}{0}  % max |local refit - stored refit| forecast
\providecommand{\rhfiLgbmInterW}{}\renewcommand{\rhfiLgbmInterW}{0.45}  % Kendall W of per-refit pair ranks, top-50 pairs
\providecommand{\rhfiLgbmInterRatio}{}\renewcommand{\rhfiLgbmInterRatio}{90}  % sum of series-pair interactions / sum main, %
\providecommand{\rhfiLgbmInterSpGlobal}{}\renewcommand{\rhfiLgbmInterSpGlobal}{0.99}
\providecommand{\rhfiLgbmInterFloor}{}\renewcommand{\rhfiLgbmInterFloor}{0.49}  % probe floor x 1e3
\providecommand{\rhfiLgbmInterAboveFloor}{}\renewcommand{\rhfiLgbmInterAboveFloor}{116}  % series pairs above the probe floor
\providecommand{\rhfiLgbmInterTopA}{}\renewcommand{\rhfiLgbmInterTopA}{HAR ladder}
\providecommand{\rhfiLgbmInterTopB}{}\renewcommand{\rhfiLgbmInterTopB}{\texttt{sumabsret}}
\providecommand{\rhfiLgbmInterTopVal}{}\renewcommand{\rhfiLgbmInterTopVal}{12.6}  % x 1e3
\providecommand{\rhfiLgbmHTop}{}\renewcommand{\rhfiLgbmHTop}{0.33}  % H^2 of the top series pair (pooled), %
\providecommand{\rhfiLgbmHMin}{}\renewcommand{\rhfiLgbmHMin}{0.02}  % min H^2 over the top-10 series pairs, %
\providecommand{\rhfiLgbmHMax}{}\renewcommand{\rhfiLgbmHMax}{0.33}  % max H^2 over the top-10 series pairs, %
\providecommand{\rhfiLgbmHMinTop}{}\renewcommand{\rhfiLgbmHMinTop}{0.19}  % top pair: interaction PD var / smaller main PD var
\providecommand{\rhfiLgbmHMinLo}{}\renewcommand{\rhfiLgbmHMinLo}{0.05}  % ... min over the top-10 series pairs
\providecommand{\rhfiLgbmHMinHi}{}\renewcommand{\rhfiLgbmHMinHi}{0.54}  % ... max
\providecommand{\rhfiLgbmHColTop}{}\renewcommand{\rhfiLgbmHColTop}{0.36}  % H^2 of the top column pair, %
\providecommand{\rhfiLgbmHMinColTop}{}\renewcommand{\rhfiLgbmHMinColTop}{0.02}  % ... min-normalised
\providecommand{\rhfiLgbmMainTopVal}{}\renewcommand{\rhfiLgbmMainTopVal}{195}  % largest global mean |phi_A| x 1e3
\providecommand{\rhfiLgbmWithinTop}{}\renewcommand{\rhfiLgbmWithinTop}{HAR ladder}
\providecommand{\rhfiLgbmWithinTopVal}{}\renewcommand{\rhfiLgbmWithinTopVal}{18.7}
\providecommand{\rhfiXgbW}{}\renewcommand{\rhfiXgbW}{0.60}  % xgb: Kendall W, per-refit series ranks
\providecommand{\rhfiXgbWCols}{}\renewcommand{\rhfiXgbWCols}{0.90}  % xgb: Kendall W, column ranks
\providecommand{\rhfiXgbTopOne}{}\renewcommand{\rhfiXgbTopOne}{HAR ladder}
\providecommand{\rhfiXgbTopOneRank}{}\renewcommand{\rhfiXgbTopOneRank}{1.0}  % mean rank
\providecommand{\rhfiXgbTopTwo}{}\renewcommand{\rhfiXgbTopTwo}{\texttt{sumret}}
\providecommand{\rhfiXgbTopTwoRank}{}\renewcommand{\rhfiXgbTopTwoRank}{3.0}  % mean rank
\providecommand{\rhfiXgbTopThree}{}\renewcommand{\rhfiXgbTopThree}{\texttt{sumabsret}}
\providecommand{\rhfiXgbTopThreeRank}{}\renewcommand{\rhfiXgbTopThreeRank}{3.9}  % mean rank
\providecommand{\rhfiXgbSpBorda}{}\renewcommand{\rhfiXgbSpBorda}{0.87}  % Spearman, mean rank vs top-10 Borda
\providecommand{\rhfiXgbSpGlobal}{}\renewcommand{\rhfiXgbSpGlobal}{0.97}  % Spearman, mean rank vs global mean |phi|
\providecommand{\rhfiXgbSpLocal}{}\renewcommand{\rhfiXgbSpLocal}{0.993}  % Spearman, mean ranks local refit vs stored
\providecommand{\rhfiXgbProbeRank}{}\renewcommand{\rhfiXgbProbeRank}{110}  % median rank of the continuous probe, 631 columns
\providecommand{\rhfiXgbProbeCols}{}\renewcommand{\rhfiXgbProbeCols}{78}  % columns above the continuous probe in most probe fits
\providecommand{\rhfiXgbProbeSeries}{}\renewcommand{\rhfiXgbProbeSeries}{35}  % exogenous series with such a column
\providecommand{\rhfiXgbTopCol}{}\renewcommand{\rhfiXgbTopCol}{\texttt{har\_ma\_1}}  % column with the lowest mean rank
\providecommand{\rhfiXgbTopColRank}{}\renewcommand{\rhfiXgbTopColRank}{1.3}
\providecommand{\rhfiXgbTopColTwo}{}\renewcommand{\rhfiXgbTopColTwo}{\texttt{har\_ma\_5}}
\providecommand{\rhfiXgbTopColTwoRank}{}\renewcommand{\rhfiXgbTopColTwoRank}{1.8}
\providecommand{\rhfiXgbPredGap}{}\renewcommand{\rhfiXgbPredGap}{0.17}  % max |local refit - stored refit| forecast
\providecommand{\rhfiXgbInterW}{}\renewcommand{\rhfiXgbInterW}{0.47}  % Kendall W of per-refit pair ranks, top-50 pairs
\providecommand{\rhfiXgbInterRatio}{}\renewcommand{\rhfiXgbInterRatio}{74}  % sum of series-pair interactions / sum main, %
\providecommand{\rhfiXgbInterSpGlobal}{}\renewcommand{\rhfiXgbInterSpGlobal}{0.99}
\providecommand{\rhfiXgbInterFloor}{}\renewcommand{\rhfiXgbInterFloor}{0.41}  % probe floor x 1e3
\providecommand{\rhfiXgbInterAboveFloor}{}\renewcommand{\rhfiXgbInterAboveFloor}{115}  % series pairs above the probe floor
\providecommand{\rhfiXgbInterTopA}{}\renewcommand{\rhfiXgbInterTopA}{HAR ladder}
\providecommand{\rhfiXgbInterTopB}{}\renewcommand{\rhfiXgbInterTopB}{\texttt{sumabsret}}
\providecommand{\rhfiXgbInterTopVal}{}\renewcommand{\rhfiXgbInterTopVal}{15.2}  % x 1e3
\providecommand{\rhfiXgbHTop}{}\renewcommand{\rhfiXgbHTop}{0.37}  % H^2 of the top series pair (pooled), %
\providecommand{\rhfiXgbHMin}{}\renewcommand{\rhfiXgbHMin}{0.02}  % min H^2 over the top-10 series pairs, %
\providecommand{\rhfiXgbHMax}{}\renewcommand{\rhfiXgbHMax}{0.68}  % max H^2 over the top-10 series pairs, %
\providecommand{\rhfiXgbHMinTop}{}\renewcommand{\rhfiXgbHMinTop}{0.39}  % top pair: interaction PD var / smaller main PD var
\providecommand{\rhfiXgbHMinLo}{}\renewcommand{\rhfiXgbHMinLo}{0.08}  % ... min over the top-10 series pairs
\providecommand{\rhfiXgbHMinHi}{}\renewcommand{\rhfiXgbHMinHi}{0.80}  % ... max
\providecommand{\rhfiXgbHColTop}{}\renewcommand{\rhfiXgbHColTop}{0.39}  % H^2 of the top column pair, %
\providecommand{\rhfiXgbHMinColTop}{}\renewcommand{\rhfiXgbHMinColTop}{0.02}  % ... min-normalised
\providecommand{\rhfiXgbMainTopVal}{}\renewcommand{\rhfiXgbMainTopVal}{192}  % largest global mean |phi_A| x 1e3
\providecommand{\rhfiXgbWithinTop}{}\renewcommand{\rhfiXgbWithinTop}{HAR ladder}
\providecommand{\rhfiXgbWithinTopVal}{}\renewcommand{\rhfiXgbWithinTopVal}{19.1}
\providecommand{\rhfiRfW}{}\renewcommand{\rhfiRfW}{0.62}  % rf: Kendall W, per-refit series ranks
\providecommand{\rhfiRfWCols}{}\renewcommand{\rhfiRfWCols}{0.93}  % rf: Kendall W, column ranks
\providecommand{\rhfiRfTopOne}{}\renewcommand{\rhfiRfTopOne}{HAR ladder}
\providecommand{\rhfiRfTopOneRank}{}\renewcommand{\rhfiRfTopOneRank}{1.0}  % mean rank
\providecommand{\rhfiRfTopTwo}{}\renewcommand{\rhfiRfTopTwo}{\texttt{sumabsret}}
\providecommand{\rhfiRfTopTwoRank}{}\renewcommand{\rhfiRfTopTwoRank}{2.6}  % mean rank
\providecommand{\rhfiRfTopThree}{}\renewcommand{\rhfiRfTopThree}{\texttt{sumret}}
\providecommand{\rhfiRfTopThreeRank}{}\renewcommand{\rhfiRfTopThreeRank}{4.0}  % mean rank
\providecommand{\rhfiRfSpBorda}{}\renewcommand{\rhfiRfSpBorda}{0.92}  % Spearman, mean rank vs top-10 Borda
\providecommand{\rhfiRfSpGlobal}{}\renewcommand{\rhfiRfSpGlobal}{0.98}  % Spearman, mean rank vs global mean |phi|
\providecommand{\rhfiRfSpLocal}{}\renewcommand{\rhfiRfSpLocal}{0.996}  % Spearman, mean ranks local refit vs stored
\providecommand{\rhfiRfProbeRank}{}\renewcommand{\rhfiRfProbeRank}{97}  % median rank of the continuous probe, 631 columns
\providecommand{\rhfiRfProbeCols}{}\renewcommand{\rhfiRfProbeCols}{71}  % columns above the continuous probe in most probe fits
\providecommand{\rhfiRfProbeSeries}{}\renewcommand{\rhfiRfProbeSeries}{29}  % exogenous series with such a column
\providecommand{\rhfiRfTopCol}{}\renewcommand{\rhfiRfTopCol}{\texttt{har\_ma\_1}}  % column with the lowest mean rank
\providecommand{\rhfiRfTopColRank}{}\renewcommand{\rhfiRfTopColRank}{1.3}
\providecommand{\rhfiRfTopColTwo}{}\renewcommand{\rhfiRfTopColTwo}{\texttt{har\_ma\_5}}
\providecommand{\rhfiRfTopColTwoRank}{}\renewcommand{\rhfiRfTopColTwoRank}{1.7}
\providecommand{\rhfiRfPredGap}{}\renewcommand{\rhfiRfPredGap}{0.11}  % max |local refit - stored refit| forecast
\providecommand{\rhfiRfInterW}{}\renewcommand{\rhfiRfInterW}{0.57}  % Kendall W of per-refit pair ranks, top-50 pairs
\providecommand{\rhfiRfInterRatio}{}\renewcommand{\rhfiRfInterRatio}{72}  % sum of series-pair interactions / sum main, %
\providecommand{\rhfiRfInterSpGlobal}{}\renewcommand{\rhfiRfInterSpGlobal}{1.00}
\providecommand{\rhfiRfInterFloor}{}\renewcommand{\rhfiRfInterFloor}{0.50}  % probe floor x 1e3
\providecommand{\rhfiRfInterAboveFloor}{}\renewcommand{\rhfiRfInterAboveFloor}{52}  % series pairs above the probe floor
\providecommand{\rhfiRfInterTopA}{}\renewcommand{\rhfiRfInterTopA}{HAR ladder}
\providecommand{\rhfiRfInterTopB}{}\renewcommand{\rhfiRfInterTopB}{\texttt{sumabsret}}
\providecommand{\rhfiRfInterTopVal}{}\renewcommand{\rhfiRfInterTopVal}{23.8}  % x 1e3
\providecommand{\rhfiRfHTop}{}\renewcommand{\rhfiRfHTop}{0.85}  % H^2 of the top series pair (pooled), %
\providecommand{\rhfiRfHMin}{}\renewcommand{\rhfiRfHMin}{0.02}  % min H^2 over the top-10 series pairs, %
\providecommand{\rhfiRfHMax}{}\renewcommand{\rhfiRfHMax}{0.85}  % max H^2 over the top-10 series pairs, %
\providecommand{\rhfiRfHMinTop}{}\renewcommand{\rhfiRfHMinTop}{1.50}  % top pair: interaction PD var / smaller main PD var
\providecommand{\rhfiRfHMinLo}{}\renewcommand{\rhfiRfHMinLo}{0.17}  % ... min over the top-10 series pairs
\providecommand{\rhfiRfHMinHi}{}\renewcommand{\rhfiRfHMinHi}{5.46}  % ... max
\providecommand{\rhfiRfHColTop}{}\renewcommand{\rhfiRfHColTop}{1.16}  % H^2 of the top column pair, %
\providecommand{\rhfiRfHMinColTop}{}\renewcommand{\rhfiRfHMinColTop}{0.06}  % ... min-normalised
\providecommand{\rhfiRfMainTopVal}{}\renewcommand{\rhfiRfMainTopVal}{224}  % largest global mean |phi_A| x 1e3
\providecommand{\rhfiRfWithinTop}{}\renewcommand{\rhfiRfWithinTop}{HAR ladder}
\providecommand{\rhfiRfWithinTopVal}{}\renewcommand{\rhfiRfWithinTopVal}{34.3}
\providecommand{\rhfiSpTreesMin}{}\renewcommand{\rhfiSpTreesMin}{0.83}  % Spearman of mean ranks, trees
\providecommand{\rhfiSpTreesMax}{}\renewcommand{\rhfiSpTreesMax}{0.96}
\providecommand{\rhfiSpTreeLinMin}{}\renewcommand{\rhfiSpTreeLinMin}{0.17}  % Spearman of mean ranks, tree vs linear
\providecommand{\rhfiSpTreeLinMax}{}\renewcommand{\rhfiSpTreeLinMax}{0.55}
\providecommand{\rhfiSpLassoRidge}{}\renewcommand{\rhfiSpLassoRidge}{0.47}
\providecommand{\rhfiHRefits}{}\renewcommand{\rhfiHRefits}{15}  % refits for H
\providecommand{\rhfiHRows}{}\renewcommand{\rhfiHRows}{100}  % evaluation points per refit
\providecommand{\rhfiHTopSeries}{}\renewcommand{\rhfiHTopSeries}{10}
\providecommand{\rhfiFitSecRf}{}\renewcommand{\rhfiFitSecRf}{212}  % mean seconds per forest refit (11 processes sharing the CPU)
\providecommand{\rhfiShapSecRf}{}\renewcommand{\rhfiShapSecRf}{182}  % mean seconds per forest refit, SHAP + interactions
\providecommand{\rhfiShapSecLgbm}{}\renewcommand{\rhfiShapSecLgbm}{0.9}  % ... LightGBM
\providecommand{\rhfiShapSecXgb}{}\renewcommand{\rhfiShapSecXgb}{0.4}  % ... XGBoost
\providecommand{\rhfiWallTreesMin}{}\renewcommand{\rhfiWallTreesMin}{133}  % wall minutes, tree stage on 11 processes
\providecommand{\rhfiWallHMin}{}\renewcommand{\rhfiWallHMin}{7}  % wall minutes, H stage
\providecommand{\rhfiWallLinMin}{}\renewcommand{\rhfiWallLinMin}{7}  % wall minutes, linear stage
\providecommand{\rhfiWorkers}{}\renewcommand{\rhfiWorkers}{11}
\fi
% 2026-10-15 (finish, user): "intraday periodicity factor" -> "diurnal adjustment" and "periodicity-adjusted" ->
% "diurnally adjusted" (term rename only; the symbol s^2_{t,c} stays); smearing mention removed.

\section{Introduction}
\label{sec:rh_intro}
% 2026-10-17 (user): prose rewritten in the style of Zhang, Zhang, Cucuringu, and Qian (2024).
\begin{itemize}
\item \textbf{Research question.} We investigate whether incorporating exogenous predictors improves intraday
  % 2026-10-17 (user, new-notes 1): parked -- the target is built from front-month E-mini futures (03_data, Realized moments). Was: "forecasts of the realized variance of the S\&P 500"
  forecasts of a futures-based (E-mini) proxy for the realized variance of the S\&P 500 over the closing half-hour interval.
  \begin{itemize}
  \item Specifically, at the 15:30 ET forecast origin on trading day $t$, we seek to predict $\RV_t$, which
    denotes the realized variance over the 15:30--16:00 closing interval (see Section~\ref{sec:rh_data}).
  \item Each intraday model is estimated across all $M = \rhdmNIntervals$ half-hour intervals of the trading
    day, and we evaluate its out-of-sample predictions on the forecast issued at 15:30 for the closing half hour.
  \item To isolate the contribution of the additional covariates, we compare the full predictor set against the
    HAR benchmark; the exogenous series constitute the exact difference between the two sets (see
    Section~\ref{sec:rh_predictors}).
  \end{itemize}
\item \textbf{Motivation for the HAR baseline.} Following \citet{Corsi2009}, the standard HAR regression models
  realized volatility solely as a linear function of its own past lags. Consequently, it omits a rich set of
  exogenous variables available by 15:30, including higher-order realized moments, liquidity measures, the
  cross-section of constituent returns, news and social-media sentiment, option-implied volatility, volatility
  demand indicators, and scheduled macroeconomic and earnings announcements.
\item \textbf{Empirical design.} For each intraday interval $(t,i)$, we consider a single forecast mapping of
  the form
  \[
    \hat y_{t,i} = f_{m,\mathcal X,\mathcal I}(\bm x_{t,i}),
  \]
  indexed by the model class $m$ and the predictor set $\mathcal X \in \{\mathrm{HAR}, \mathrm{Full}\}$.
  % 2026-10-16 (user, E12): parked -- $\hat y_t = f_{m,\mathcal X,\mathcal I}(\bm x_t)$.
  \begin{itemize}
  \item The primary empirical results reported in Section~\ref{sec:rh_results} focus on the pooled intraday
    estimation scheme, $\mathcal I = \{1, \dots, M\}$.
  % 2026-10-18 (user, new-notes 10): parked -- closing cadence only in the closing section. Was: "\item As our primary benchmark forecast, we employ the daily ordinary least squares (OLS) regression on the"
  \item As our primary benchmark forecast, we employ the ordinary least squares (OLS) regression on the
    % 2026-10-17 (user): parked -- closing-interval material (relabeled: closing-interval HAR benchmark). Was: "estimated exclusively on the closing interval (see"
    HAR benchmark predictor set, estimated exclusively on the closing interval (the closing-interval HAR benchmark; see
    Sections~\ref{sec:rh_methods_models} and~\ref{sec:rh_results_eval}).
  \item We evaluate predictive accuracy via the average QLIKE loss over the out-of-sample evaluation span
  % 2026-10-20 (user): parked -- Table 1 scored from 2012 (TABLE1_SWITCH.md, approved). Was: "$\mathcal T_{\mathrm{oos}}$ from \rhoosSpanStart{} to \rhoosSpanEnd{} ($T_{\mathrm{oos}} = \rhoosDays$ trading days; see Section~\ref{sec:rh_results_eval})."
    $\mathcal T_{\mathrm{oos}}$ from \rhoosSpanStart{} to \rhoosSpanEnd{} ($T_{\mathrm{oos}} = \rhoosDays$
    % 2026-10-22 (user, E36): parked -- "the accuracy of the pooled linear models in Table~\ref{tab:rh_oos}"; Table 1 now holds the tree ensembles too.
    trading days; see Section~\ref{sec:rh_results_eval}); the accuracy of the pooled linear models and tree
    ensembles in Table~\ref{tab:rh_oos} is scored on every regular-hours bar from \rhorSpanStart{} to \rhorSpanEnd{}
    (\rhorDays{} trading days, a rolling window of \rhorWindowDays{} trading days).
  \end{itemize}
\item \textbf{Main findings.}
  \begin{itemize}
  % 2026-10-22 (user, E34: one design for the 12 pooled forecasts): parked -- "across every intraday model class" and "the elastic net achieves a statistically significant reduction" no longer hold with the pooled linear models on the trees' design (make_rehaul_results.py, E34 CLAIM); counts printed. Was: "\item Incorporating the full predictor set reduces the average out-of-sample QLIKE loss relative to the HAR benchmark across every intraday model class (\rhrAfPoolPairs{} classes in total, with \rhrAfPoolLowerSig{} of these reductions statistically significant at the 5\% level; see Section~\ref{sec:rh_results}). Among the linear specifications, the elastic net achieves a statistically significant reduction; among the tree ensembles, the reduction is statistically significant at the 5\% level in \rhrAfPoolTreeLowerSig{} of the three classes."
  \item Incorporating the full predictor set reduces the average out-of-sample QLIKE loss relative to the HAR
    benchmark in \rhrAfPoolLower{} of the \rhrAfPoolPairs{} intraday model classes, with \rhrAfPoolLowerSig{}
    of these reductions statistically significant at the 5\% level (see Section~\ref{sec:rh_results}).
    % 2026-10-22 (user, E34): parked -- with the base-5 pooled trees their reductions are no longer the largest or significant (make_rehaul_results.py, E34 CLAIM); no ranking. Was: "The tree-based ensembles exhibit the largest relative improvements, while among the linear specifications, the elastic net achieves a statistically significant reduction."
    The reduction is statistically significant at the 5\% level in \rhrAfPoolLinLowerSig{} of the three linear
    classes and in \rhrAfPoolTreeLowerSig{} of the three tree classes.
  % 2026-10-19 (user): parked -- the MSE finding, with the QLIKE/squared-error divergence bullet of 06_results.tex
  % and the abstract's MSE clause. Was:
  % \item In contrast, evaluation under mean squared error (MSE) does not corroborate these gains: the full
  %   predictor set increases the average MSE in \rhrAfPoolMseHigher{} of the \rhrAfPoolPairs{} model classes,
  %   and none of the MSE differences is statistically significant.
  % 2026-10-17 (user): parked -- closing-interval material (relabeled: closing-interval HAR benchmark). Was: "than the closing-interval benchmark forecast."
  \item Every pooled intraday forecast incurs a higher average QLIKE loss than the closing-interval HAR
    % 2026-10-22 (user, E34): parked -- with the base-5 pooled trees the gap on the full predictor set is significant for \rhrVsBmPoolTreeFullSig{} of the three tree ensembles, not all (make_rehaul_results.py, E34 CLAIM). Was: "benchmark. While this performance gap is statistically insignificant at the 5\% level for the linear models, it is significant for the tree ensembles. Notably, augmenting the intraday models with the full predictor"
    % 2026-10-22 (user, E34: one design for the 12 pooled forecasts): parked -- "insignificant ... for the linear models" and "narrows this gap across all model classes" no longer hold (make_rehaul_results.py, E34 CLAIM); counts printed. Was: "benchmark. Under the full predictor set, this performance gap is statistically insignificant at the 5\% level for the linear models and significant for \rhrVsBmPoolTreeFullSig{} of the three tree ensembles. Notably, augmenting the intraday models with the full predictor set narrows this gap across all model classes (see Section~\ref{sec:rh_results_classes})."
    benchmark. Under the full predictor set, this performance gap is statistically significant at the 5\% level
    for \rhrVsBmPoolLinFullSig{} of the three linear models and for \rhrVsBmPoolTreeFullSig{} of the three tree
    ensembles. Augmenting the intraday models with the full predictor set narrows this gap in
    \rhrVsBmPoolNarrow{} of the \rhrAfPoolPairs{} model classes (see Section~\ref{sec:rh_results_classes}).
  \end{itemize}
\item \textbf{Economic application.} To assess the economic value of the forecasts on trading days with a
  % 2026-10-18 (user, second-list 20): parked -- SPXW is the weekly PM-settled series the strategy trades; SPX is the monthly. Was: "same-day SPX expiration"
  same-day SPXW expiration, we construct a trading strategy that takes a long or short position in a
  zero-days-to-expiration (0DTE) at-the-money straddle based on the sign of the difference between the variance
  forecast and option-implied variance (see Section~\ref{sec:rh_econ}). Using the full predictor set increases
  the annualized Sharpe ratio of this strategy relative to the HAR benchmark in every intraday model class (see
  Section~\ref{sec:rh_econ_results}).
% 2026-10-22 (user, E34: one design for the 12 pooled forecasts): parked -- the note named one class. Was: "\editc{orange}{N15c}{Problem: ``every intraday model class'' no longer holds under the per-year Sharpe ratio (pooled random forest $\Delta\SR = \rhrAfRfPoolDSh$); see note N15c in Section~\ref{sec:rh_econ}.}"
% 2026-10-22 (user, D6): parked -- the note under the per-year ratio D1; under the Sharpe ratio of note D6 the sentence holds in point estimates. Was: "\editc{orange}{N15c}{Problem: ``every intraday model class'' does not hold: the full predictor set raises the pooled Sharpe ratio in \rhrAfPoolShHigher{} of the \rhrAfPoolPairs{} classes ($\Delta\SR$ from \rhrAfPoolDShMin{} to \rhrAfPoolDShMax); see note N15c in Section~\ref{sec:rh_econ}.}"
\editc{orange}{N15c}{Problem: under the Sharpe ratio of note D6, the full predictor set raises the pooled Sharpe ratio in \rhrAfPoolShHigher{} of the \rhrAfPoolPairs{} classes ($\Delta\SR$ from \rhrAfPoolDShMin{} to \rhrAfPoolDShMax), so ``every intraday model class'' holds in point estimates; \rhrAfPoolShHigherSig{} of these gains are statistically significant at the 5\% level (bootstrap test), and under the previous ratio the count was \rheDefDOnePos{} of \rheDefN{}. See notes N15c and D6 in Section~\ref{sec:rh_econ}.}
\item \textbf{Paper organization.} The remainder of the paper is organized as follows.
  Section~\ref{sec:rh_related} reviews the related literature. Section~\ref{sec:rh_data} describes the data and
  target construction, and Section~\ref{sec:rh_predictors} details the benchmark and exogenous predictors.
  Section~\ref{sec:rh_methods} presents the forecasting models and rolling estimation protocol.
  % 2026-10-17 (user): parked -- closing-interval material (the variable-importance analysis covers the closing-interval models and is parked). Was: "reports the out-of-sample forecasting results and variable-importance analysis."
  Section~\ref{sec:rh_results} reports the out-of-sample forecasting results.
  Section~\ref{sec:rh_econ} evaluates the economic value of the forecasts in the 0DTE straddle strategy, and
  Section~\ref{sec:rh_conclusion} concludes.
\end{itemize}
% 2026-10-16 (user): the introduction is about exogenous features in the intraday forecast,
% not the 24-hour panel and not the straddle as the question. Parked verbatim below.
% \begin{itemize}
% \item \textbf{Question.} At the forecast origin, 15:30 ET on trading day $t$, we forecast the realized variance of the
%   S\&P 500 over the closing half hour (15:30--16:00 ET), $\E_t[\RV_t]$, where $\E_t[\cdot] \equiv \E[\,\cdot \mid
%   \mathcal F_t]$ and the information set $\mathcal F_t$ holds all data through 15:30 on day $t$.
%   \begin{itemize}
%   % 2026-10-15 (consolidation): parked -- the upper limit 30 contradicts Data and Methods (N_{t,i} <= 30 varies).
%   % \item $\RV_t = \sum_{k=1}^{30} r_{t,c,k}^2$: the sum of squared one-minute log returns $r_{t,c,k}$ of E-mini
%   %   futures over the closing interval $i = c$ \citep{AndersenBollerslevDieboldLabys2003}.
%   % 2026-10-16 (formula pass): parked -- the sum repeats Data. Was:
%   % RV_t = sum_{k=1}^{N_{t,c}} r_{t,c,k}^2, N_{t,c} <= 30 (Andersen, Bollerslev, Diebold and Labys 2003).
%   \item $\RV_t$: the realized variance of the closing interval (Section~\ref{sec:rh_data}).
%   % 2026-10-15 (notation pass): parked -- a straddle's payoff is |S_T - K|, set by the closing half hour's move,
%   % not by its realized variance.
%   % \item A 0DTE SPXW straddle opened at 15:30 settles on exactly that variance.
%   % 2026-10-15 (consolidation): parked -- "delta-neutral" holds only approximately and only at entry (07_economics).
%   % \item A zero-days-to-expiration (0DTE) SPX Weeklys (SPXW) straddle position, delta-neutral and at the money,
%   %   entered at 15:30 and held to expiration, settles at the 16:00 close: its payoff is set by the index move over
%   %   exactly this half hour.
%   \item A zero-days-to-expiration (0DTE) SPX Weeklys (SPXW) at-the-money straddle position, approximately
%     delta-neutral at entry, entered at 15:30 and held unhedged to expiration, settles at the 16:00 close: its payoff
%     is set by the index move over exactly this half hour (Section~\ref{sec:rh_econ_strategy}).
%   % 2026-10-16 (formula pass): parked -- the display defined hat sigma^2_t, then q_t from hat sigma^2_t, then
%   % the next bullet defined y_t and s^2_{t,c}, so the position took two further substitutions. The forecast and
%   % the adjustment are defined once in Methods; the sign rule is defined once in Economic value. Was:
%   % hat sigma^2_t = hat y_t^2 s^2_{t,c}, q_t = sign(hat sigma^2_t - IV_t^2);
%   % y_t = sqrt(RV_t / s^2_{t,c}); s^2_{t,c} the mean of RV over the previous 20 trading days.
%   \item The variance forecast, the diurnal adjustment and the target are defined in
%     Section~\ref{sec:rh_methods_setup}. The position is the sign strategy of
%     Section~\ref{sec:rh_econ_strategy}: $q_t = +1$ buys the straddle, $q_t = -1$ sells it ($q_t = -1$ on ties),
%     and the daily return is $q_t R_t$, with $R_t$ the straddle's simple return from 15:30 to settlement.
%   \item Sample: $T_{\mathrm{oos}} = \cmNdays$ out-of-sample trading days with an SPXW expiration, \rheSampleStart{}
%     to \rheSampleEnd, early closes excluded.
%   \end{itemize}
% \item \textbf{Why it matters.}
%   \begin{itemize}
%   \item Implied variance exceeds the forecast, $\hat\sigma^2_t < \mathrm{IV}_t^2$, on
%     \cmSellPaperMin--\cmSellPaperMax\% of trading days, depending on the model, in line with a positive variance
%     risk premium; a forecast adds value mainly by identifying the days on which to buy.
%   \item 0DTE options account for a large share of S\&P 500 index option volume
%     \citep{BandiFusariReno2023, BeckmeyerBrangerGayda2023}.
%     % 2026-10-15 (bib pass): parked -- the typed citation, now \citep keys in references.bib. Was:
%     % cite: Bandi, Fusari and Reno (2023); Beckmeyer, Branger and Gayda (2023) -- not yet in references.bib.
%   \item The payoff depends on a single, pre-specified window, so the forecast target and the holding period
%     coincide.
%   \end{itemize}
% \item \textbf{Two specification choices, one forecast function:}
%   \[
%     \hat y_t = f_{m,\mathcal X,\mathcal I}(\bm x_t), \qquad \bm x_t \in \mathbb{R}^p .
%   \]
%   \begin{itemize}
%   \item Predictor set $\mathcal X \in \{\text{HAR}, \text{Full}\}$: the HAR benchmark, $p = \cmPNewBase$ (HAR
%     components and calendar dummies), vs the full predictor set, $p = \cmPNewAll$ (adding \rhfiNExog{} exogenous
%     series in \rhfiNBuckets{} groups).
%   \item Estimation sample $\mathcal I \in \bigl\{\{c\}, \{1, \dots, M\}\bigr\}$: the closing-interval model (the
%     closing half hour only, one observation per trading day) vs the intraday model (every 30-minute
%     interval, $M = 48$).
%     % 2026-10-15 (consolidation): M = 48 is typed; no macro yet (generator: \rhdmNIntervals in make_rehaul_data_methods.py).
%   \item Model class $m$: linear (OLS, ridge, lasso, elastic net), tree ensemble (random forest, LightGBM, XGBoost)
%     or neural network (MinGRU), crossed with $(\mathcal X, \mathcal I)$.
%   \end{itemize}
% % 2026-10-16 (formula pass): parked -- the display defined the average, then L_t, then the QLIKE function, and
% % the function is defined once in Related work, the scored loss once in Methods, and the average once in
% % Forecast evaluation. Was:
% % overline{QLIKE} = (1/T_oos) sum L_t, L_t = QLIKE(RV_t, hat sigma^2_t),
% % QLIKE(RV, hat sigma^2) = RV/hat sigma^2 - log(RV/hat sigma^2) - 1.
% \item \textbf{Evaluation loss} \citep{Patton2011}: average QLIKE over the evaluation span
%   $\mathcal T_{\mathrm{oos}}$ (Sections~\ref{sec:rh_methods_setup} and~\ref{sec:rh_results_eval}).
%   \begin{itemize}
%   \item Economic evaluation: the annualized Sharpe ratio $\SR$ of $q_t R_t$ at mid-quotes.
%   \end{itemize}
% \item \textbf{Contributions.}
%   % 2026-10-15 (notation pass): parked -- recast as question, data, method, findings and economic value; each old
%   % item is kept in substance below (indicators -> Data; back-transform, causal tuning test, MinGRU -> Method).
%   % \item Value, availability and activity ladders: 47 exogenous series in one $x_t$.
%   % \item Plain back-transform $\hat y_t^2 B_t$: no smearing, no recalibration.
%   % \item Causal walk-forward test of tuning: refits, re-tuning, search budgets.
%   % \item Native GPU MinGRU, regularised like the linear stack.
%   \begin{itemize}
%   \item \emph{Question.} We forecast the variance of the closing half hour, the risk a 0DTE SPX straddle opened at
%     15:30 carries to settlement.
%   \item \emph{Data.} One predictor vector $\bm x_t$ holds, for each of \rhfiNExog{} exogenous series, HAR
%     components of its value, of its missing-data indicator and, for zero-inflated series, of an activity
%     indicator; $T_{\mathrm{oos}} = \cmNdays$ out-of-sample trading days.
%   \item \emph{Method.} A rolling-window out-of-sample exercise, re-estimated daily, in which every hyperparameter is
%     selected on data before the forecast origin. We measure what re-estimation frequency, hyperparameter
%     re-selection and search budgets add out of sample. Forecasts are retransformed as $\hat y_t^{\,2} s^2_{t,c}$,
%     with no recalibration. The model classes include a MinGRU recurrent network \citep{Feng2024},
%     regularized like the linear models.
%     % 2026-10-15 (user): parked -- no smearing mentions. Was: "... as $\hat y_t^{\,2} s^2_{t,c}$, with no smearing or
%     % recalibration."
%     % 2026-10-15 (bib pass): parked -- the typed citation, now a \citep key in references.bib. Was:
%     % MinGRU: Feng et al. (2024) -- not yet in references.bib.
%   \item \emph{Findings.} The full predictor set lowers $\overline{\QLIKE}$ relative to the HAR benchmark in
%     \fqsQLower{} of \fqsPairs{} pairs, \fqsQLowerSig{} of them significantly at the 5\% level. The closing-interval
%     model is more accurate than the intraday model; for ridge on the HAR benchmark, $\overline{\QLIKE}$
%     changes by \cmLdPerBarBaseQ\% ($\DM = \cmLdPerBarBaseDM$).
%   \item \emph{Economic value.} \cmMasterVsShortAbove{} of the \cmMasterVsShortN{} forecasts earn a higher $\SR$
%     than the unconditional short-straddle benchmark ($q_t \equiv -1$), \cmMasterVsShortAboveSig{} of them
%     significantly at the 5\% level. Accuracy and economic value are only weakly related across forecasts.
%   \end{itemize}
% \item \textbf{Predictor set: full predictor set vs HAR benchmark} (each pair differs only in $\mathcal X$).
%   \begin{itemize}
%   \item $\Delta\overline{\QLIKE} = \overline{\QLIKE}_{\text{Full}} - \overline{\QLIKE}_{\text{HAR}} < 0$ in
%     \fqsQLower{} of \fqsPairs{} pairs; median relative change \fqsMedDQPct\%.
%   \item \fqsQLowerSig{} of these reductions are significant at the 5\% level (Diebold--Mariano statistic with the
%     Harvey--Leybourne--Newbold correction); \fqsQHigherSig{} of the \fqsQHigher{} increases are.
%   \item $\Delta\SR > 0$ in \fqsShHigher{} of \fqsPairs{} pairs; median \fqsMedDSh; \fqsShHigherSig{} significant at
%     the 5\% level (HAC $t$-statistic on the daily return differential).
%   \item All \cmMasterVsShortAboveSig{} forecasts that significantly beat the short-straddle benchmark
%     ($q_t \equiv -1$) use the full predictor set.
%   \end{itemize}
% \item \textbf{Estimation sample: closing-interval vs intraday model.}
%   \begin{itemize}
%   \item The closing-interval model has the lower $\overline{\QLIKE}$ wherever both are estimated.
%   \item The difference is significant at the 5\% level on the HAR benchmark in every model class; for ridge,
%     $\overline{\QLIKE}$ changes by \cmLdPerBarBaseQ\% ($\DM = \cmLdPerBarBaseDM$).
%   \item The difference in $\SR$ is not statistically significant.
%   \end{itemize}
% \item \textbf{Statistical accuracy and economic value are weakly related.}
%   \begin{itemize}
%   \item Spearman rank correlation of $\overline{\QLIKE}$ and $\SR$ across the \cmMasterN{} forecasts:
%     \cmMasterRankCorr{} (negative: a lower loss goes with a higher $\SR$).
%   \item Of the \cmLdN{} single specification changes, \cmLdResolved{} moves $\SR$ significantly.
%   \end{itemize}
% \item \textbf{Roadmap:} data (Section~\ref{sec:rh_data}), predictors (Section~\ref{sec:rh_predictors}), methods
%   (Section~\ref{sec:rh_methods}), statistical results, first for the predictor set and then for the estimation
%   sample (Section~\ref{sec:rh_results}), the economic value of the straddle strategy (Section~\ref{sec:rh_econ}),
%   variable importance (Section~\ref{sec:rh_fi}) and the conclusion (Section~\ref{sec:rh_conclusion}).
% \end{itemize}

